\subsection{A correction from a second divided form}
\label{ssec:dividedcorrection}

The same finite eigenstate supplies an error predictor for every simple
nonlattice transform root. The first zero is a special case; uniform
accuracy in a growing root index requires an additional remainder estimate.
The following construction uses finite operator data and no reference-zero
errors or fitted conversion constant.

Write $L=\log\lambda^2$, $\alpha_j=2\pi j/L$, and let $a$ be a real even
unit eigenvector of the full finite matrix $A$, with $Aa=\varepsilon_N a$.
Set
\[
 F(z)=\frac{2}{\sqrt L}\sin(zL/2)
       \sum_{|j|\le N}\frac{a_j}{z-\alpha_j}.
\]
At a simple nonlattice root $\nu>0$ of $F$, define
\begin{equation}\label{eq:divided-observables}
 b_j=\frac{a_j}{\alpha_j-\nu},\quad
 S=a^Tb,\quad J=b^TAb,\quad f=F'(\nu),\quad
 h=\frac{F''(\nu)}{F'(\nu)},\quad \kappa=1-\nu h.
\end{equation}
For $\kappa\ne0$, the corrected signed displacement estimate is
\begin{equation}\label{eq:divided-predictor}
 \boxed{\widehat\delta=
   \frac{\nu(J-f^2)-\varepsilon_N S}{\kappa f^2}.}
\end{equation}
Its relative-error estimate is $|\widehat\delta|/\nu$; a known zeta zero
enters only the subsequent scoring. The bare local estimate is
$-\varepsilon_NS/f^2$. All terms in~\eqref{eq:divided-predictor} must use
the same normalization, and the subtraction $J-f^2$ needs sufficient
source and working precision.

\begin{proposition}[Finite corrected-displacement identity]
\label{prop:divided-correction}
Let $\gamma>0$ be a correctly matched simple real zeta ordinate and put
$d=\gamma-\nu$ and $B=F(\gamma)/d$, with $B=f$ at $d=0$.
Over the full multiset $z_\rho=(\rho-1/2)/i$ of nontrivial zeta zeros,
define
\[
 R(\nu)=\sum_{z_\rho\ne\pm\gamma}
         \frac{F(z_\rho)^2}{z_\rho-\nu},\qquad
 R'(\nu)=\sum_{z_\rho\ne\pm\gamma}
         \frac{F(z_\rho)^2}{(z_\rho-\nu)^2}.
\]
The prime differentiates the rational weights with $F$ and the excluded
pair fixed. Set
\begin{align}
 E_{\rm loc}={}&\frac{2\nu dB^2}{2\nu+d}
   -\nu\left\{B^2\left(1+\frac{d^2}{(2\nu+d)^2}\right)-f^2\right\}
   -\kappa f^2d.\label{eq:divided-local}
\end{align}
Then, without assuming RH,
\begin{equation}\label{eq:divided-remainder}
 \widehat\delta-(\nu-\gamma)
 =\frac{\nu R'(\nu)-R(\nu)-E_{\rm loc}}{\kappa f^2}.
\end{equation}
For fixed nonzero $f,\nu$, one has $E_{\rm loc}=O(d^2)$.
\end{proposition}
\begin{proof}
The secular root condition gives
$F_b(z)=F(z)/(z-\nu)$ by partial fractions. The general Weil Hermitian
identity~\cite[Section~1.2]{AlpogeFurman2026}, applied to real-type transforms,
therefore gives
\[
 \varepsilon_NS=\frac{2\nu d B^2}{2\nu+d}+R(\nu),\qquad
 J=B^2\left(1+\frac{d^2}{(2\nu+d)^2}\right)+R'(\nu).
\]
Both sums converge absolutely for these finite Fourier transforms. At an
off-line zero the summand is a square, not an absolute square; conjugate
pairs make the total real. The analytic divided transform supplies the
displayed value at coincidence. Substitution gives
\eqref{eq:divided-remainder}. Finally,
$B=f+\tfrac12F''(\nu)d+O(d^2)$, so the constant and linear terms in
$E_{\rm loc}$ cancel.
\end{proof}

\paragraph{Computing the second form.}
The rank-two commutator, with $D=\operatorname{diag}(\alpha_j)$,
is $[D,A]=g\mathbf1^T-\mathbf1g^T$. Its generator can be recovered from
one column: $g_j=\alpha_jA_{j0}$, $g_0=0$. Since $\mathbf1^Tb=0$,
\begin{equation}\label{eq:divided-fastform}
 J=\varepsilon_N\|b\|^2+(g^Tb)r'(\nu),\qquad
 r(z)=\sum_j\frac{a_j}{z-\alpha_j}.
\end{equation}
Indeed $(D-\nu I)Ab=\varepsilon_Na-\mathbf1(g^Tb)$; multiply by
$(D-\nu I)^{-1}$ and pair with $b$. Thus evaluation costs $O(N)$ per
root once the generator and eigenstate are available. The full matrix is
still needed to determine that eigenstate: the commutator does not fix
the diagonal or the eigenvalue decay.

\paragraph{Why the correction helps, and its limit.}
Under RH, put $E_m=2F(\gamma_m)^2$. Pairing remote zeros yields
\begin{equation}\label{eq:divided-secondweight}
 \nu R'-R=2\nu^3\sum_{m\ne k}
       \frac{E_m}{(\gamma_m^2-\nu^2)^2}\ge0.
\end{equation}
The contribution with $\gamma_m>H>\nu$ is at most
$2\nu^3E_{>H}/(H^2-\nu^2)^2$. The bare estimate instead has a
remote bound $\nu E_{>H}/(H^2-\nu^2)$ before division by $f^2$.
The correction cancels the leading behavior $R(\nu)\approx\nu M_2$ of
a distant spectral distribution. An unconditional bound needs the
absolute size of the signed-complex sum, and finite positive matrices
do not supply that bound. The identities above are unconditional; smallness
of their remainder uniformly in $\lambda^2,N,k$ remains open.
The denominator $\kappa$, local curvature, correct root matching and
numerical assembly error must also be controlled.

\paragraph{Retained-data test, 10 September 2026.}
The first twenty converged stored roots in each of five configurations
give Table~\ref{tab:divided-test}. All 100 signs agree. The reported
errors compare error magnitudes, not percentages of large digit depths.
The source eigenvectors, root records and matrix ingredients are linked
by their artifact identities; the postprocessing was repeated with 100
fewer decimal digits, giving identical exported 32-significant-digit
prediction results. This is point arithmetic on common retained sources,
not independent assembly or interval certification.

\begin{table}[htbp]
\centering\small
\begin{tabular}{@{}rrrr@{}}
\toprule
$\lambda^2$ & $N$ & source bits & maximum relative prediction error \\
\midrule
13 & 120 & 3386 & $0.348670\%$ \\
100 & 500 & 3386 & $0.000275\%$ \\
1000 & 800 & 6708 & $0.000061\%$ \\
1000 & 890 & 6708 & $0.000047\%$ \\
1200 & 970 & 6708 & $0.000035\%$ \\
\bottomrule
\end{tabular}
\caption{Unfitted corrected error estimates, first twenty converged roots
per configuration. Reference ordinates enter scoring, not the predictor.}
\label{tab:divided-test}
\end{table}

At $(\lambda^2,N)=(1000,890)$, selected higher-root diagnostics reduce
the error-estimate discrepancy from $46.10\%$ to $0.00790\%$ at $k=500$.
It is $0.03039\%$ at $k=575$, but rises to $1.41441\%$ at $k=600$ and
$2.23735\%$ at $k=610$. Those last two fail a one-percent accuracy gate.
These higher roots use reference-seeded Newton refinement of the retained
transform, not independent discovery or a root-counting certificate; the
selected tests do not prove every intervening index passes. The result
advances the root-conversion channel. A formula predicting the eigenvalue
and eigenvector from cutoff and resolution alone is still required.
